\section{Conclusion}\label{sec:conclusion}

A reserve-currency sovereign does not run into a debt wall. It runs into a choice about who bears a fiscal gap, and the monetary response decides by whom (Table~\ref{tab:whopays}). The table does not rank the routes: which is best depends on the distortions of each and on who should bear the cost, and these are political judgments. What the paper adds is the prices, which make the choice explicit rather than implicit.

\begin{table}[htbp]
\centering
\caption{Who bears a fiscal gap}\label{tab:whopays}
\footnotesize
\begin{tabularx}{\textwidth}{>{\raggedright\arraybackslash}X>{\raggedright\arraybackslash}X>{\raggedright\arraybackslash}X>{\raggedright\arraybackslash}X}
\toprule
Route & Who pays & Decided by & Price or size in this paper \\
\midrule
Fiscal adjustment & Taxpayers or program beneficiaries & Legislation & Threshold $\varphi^*$: about 0.47 of marginal interest cost (United States, 2023–25) \\
Inflation, tolerated by the central bank & Holders of nominal claims: savers, bondholders, currency holders (part of them abroad) & Central bank, without a legislative vote & About 0.2 pp of inflation a year for a decade per 0.1 of missing offset; Japan 2022–25: about 18\% of GDP from holders of nominal claims \\
Higher real rates, when the central bank tightens & Taxpayers, through interest paid to bondholders and, on reserves, to banks; workers, through slower growth & Central bank & Above $\kappa^*$ = 1/R (0.46 in the United States in 2025), inflation cannot close the gap at all; United States and UK 2021–25: most of a 10–14\% of GDP transfer returned; UK interest on reserves alone 5.7\% of GDP \\
\bottomrule
\end{tabularx}
\end{table}


Three findings organize these prices. The fiscal layer is maturity-free: the adjustment needed depends on $r-g$, on how rates respond to debt and on the debt level, and in the United States in 2023--25 it is ordinary by the standard of 1980--2007. The price of the inflation route is set by the consolidated balance sheet. QE funded the Fed's book with interest-bearing reserves rather than currency, which raised that price to its highest level since at least 1980 and lowered, by about two fifths, the monetary tightening the route can survive; over ten years or more, the threshold is now below a Taylor rule. The UK and Japan show the same mechanism, whatever their debt management. And which route is taken depends on the central bank. Where it raised real rates after 2021, the transfer from the inflation was largely returned, mostly through interest on reserves; where yields were held down, holders of nominal claims bore it. The durable resolutions are a fiscal response that meets the threshold, or a one-time surprise revaluation of the debt.

Several issues remain open. The U.S. level of the fiscal response is not identified, which is why results are reported as a function of it. The debt sensitivity of rates may differ across regimes, and ranking today's requirement first needs today's $\psi$ to be at least about 2bp. The prices are first order and undiscounted and hold currency demand fixed, and the balance-sheet counterfactuals ignore term-premium effects. Extending the Japanese series before 2020 and adding the UK's Treasury bills are natural next steps, and three sovereigns are not yet the panel that could identify the fiscal response across countries.
